\documentclass[11pt,a4paper]{article}
\newcommand{\fichetype}{Corrigé détaillé}
% Thème A4 commun, librement modifiable.
\usepackage[T1]{fontenc}
\usepackage[utf8]{inputenc}
\usepackage[french]{babel}
\usepackage{lmodern,amsmath,amssymb,mathtools,tikz,pgfplots}
\usepackage[margin=2.2cm,headheight=15pt]{geometry}
\usepackage{xcolor,fancyhdr,enumitem,hyperref,needspace}
\usetikzlibrary{arrows.meta,positioning,calc}
\pgfplotsset{compat=1.18,width=11cm,height=6cm}
\definecolor{Ink}{HTML}{202124}
\definecolor{Blue}{HTML}{356B9E}
\definecolor{Muted}{HTML}{737579}
\color{Ink}
\renewcommand{\familydefault}{\sfdefault}
\setlength{\parindent}{0pt}
\setlength{\parskip}{6pt}
\setlist{itemsep=3pt,topsep=4pt}
\pagestyle{fancy}\fancyhf{}
\fancyhead[L]{\small\color{Muted}MATHS LYCÉE}
\fancyhead[R]{\small\color{Muted}2026--2027}
\fancyfoot[L]{\small\color{Muted}\fichetype}
\fancyfoot[R]{\small\thepage}
\renewcommand{\headrulewidth}{0.3pt}
\hypersetup{colorlinks=true,urlcolor=Blue,linkcolor=Blue}
\newcommand{\R}{\mathbb{R}}
\newcommand{\N}{\mathbb{N}}
\newcommand{\E}{\mathbb{E}}
\newcommand{\V}{\mathrm{V}}
\newcommand{\exercice}[3]{\Needspace{6\baselineskip}\vspace{0.5cm}{\color{Blue}\small\textbf{EXERCICE #1}\hfill #2}\par{\large\bfseries #3}\par}

\begin{document}
{\small\color{Blue}PREMIÈRE SPÉCIALITÉ}\par{\LARGE\bfseries Calculer une dérivée}\par{\large Corrigé détaillé}\par\vspace{0.2cm}\hrule\vspace{0.4cm}
\par\noindent\begin{minipage}{\linewidth}\exercice{1}{Application}{Références}
{\small\color{Muted}\textbf{Énoncé.} Dériver $4x^3-5x+2$, $2/x$ et $3\sqrt{x}$. Préciser les domaines.}
\par\medskip
\textbf{Résolution.}\par
Le polynôme est dérivable sur $\R$ et sa dérivée vaut $12x^2-5$. Sur $\R\setminus\{0\}$, $(2/x)\prime=-2/x^2$. Enfin $3\sqrt{x}$ est définie sur $[0,+\infty[$, dérivable sur $]0,+\infty[$ et sa dérivée vaut $3/(2\sqrt{x})$.
\end{minipage}\par
\par\noindent\begin{minipage}{\linewidth}\exercice{2}{Application}{Un produit}
{\small\color{Muted}\textbf{Énoncé.} Dériver $(x^2-1)(x+4)$ de deux façons.}
\par\medskip
\textbf{Résolution.}\par
Par produit : $2x(x+4)+(x^2-1)=3x^2+8x-1$. Par développement, $f=x^3+4x^2-x-4$, d’où la même dérivée. Les deux calculs sont valables sur $\R$.
\end{minipage}\par
\par\noindent\begin{minipage}{\linewidth}\exercice{3}{Entraînement}{Quotient}
{\small\color{Muted}\textbf{Énoncé.} Dériver $(3x-2)/(x+1)$ et étudier le signe de la dérivée.}
\par\medskip
\textbf{Résolution.}\par
Le domaine exclut $-1$. Le numérateur de la dérivée vaut $3(x+1)-(3x-2)=5$. Donc $f\prime=5/(x+1)^2>0$ pour $x\ne-1$. La positivité vaut séparément sur les deux intervalles du domaine.
\end{minipage}\par
\par\noindent\begin{minipage}{\linewidth}\exercice{4}{Entraînement}{Racine composée}
{\small\color{Muted}\textbf{Énoncé.} Donner les domaines de définition et dérivabilité de $\sqrt{5-2x}$, puis sa dérivée.}
\par\medskip
\textbf{Résolution.}\par
$5-2x\ge0$ équivaut à $x\le5/2$. Pour dériver la racine, il faut $5-2x>0$, donc $x<5/2$. Le facteur intérieur est $-2$, d’où $f\prime=-1/\sqrt{5-2x}$. Le bord $5/2$ ne fournit pas de nombre dérivé fini.
\end{minipage}\par
\par\noindent\begin{minipage}{\linewidth}\exercice{5}{Entraînement}{Repérer une erreur}
{\small\color{Muted}\textbf{Énoncé.} Un élève écrit $((x+1)/x)\prime=1$. Corriger et vérifier par simplification.}
\par\medskip
\textbf{Résolution.}\par
La règle du quotient donne $[x-(x+1)]/x^2=-1/x^2$. On retrouve ce résultat en écrivant $(x+1)/x=1+1/x$. L’élève a divisé les dérivées, ce qui n’est pas une règle de dérivation.
\end{minipage}\par
\par\noindent\begin{minipage}{\linewidth}\exercice{6}{Entraînement}{Tangente}
{\small\color{Muted}\textbf{Énoncé.} Pour $f(x)=1/(x^2+1)$, déterminer la tangente en $1$.}
\par\medskip
\textbf{Résolution.}\par
$f(1)=1/2$. La règle de l’inverse donne $f\prime(x)=-2x/(x^2+1)^2$, donc $f\prime(1)=-1/2$. La tangente est $y=1/2-(x-1)/2=1-x/2$.
\end{minipage}\par
\par\noindent\begin{minipage}{\linewidth}\exercice{7}{Entraînement}{Contrôle numérique}
{\small\color{Muted}\textbf{Énoncé.} Pour $f(x)=x^3$, calculer le taux entre $2$ et $2+h$, puis comparer à la dérivée pour $h=0,01$.}
\par\medskip
\textbf{Résolution.}\par
En développant $(2+h)^3$, le taux vaut $12+6h+h^2$. Pour $h=0,01$, il vaut $12,0601$. La dérivée en $2$ vaut $3\times2^2=12$. Le petit écart est attendu : le taux moyen n’est pas exactement la dérivée.
\end{minipage}\par
\par\noindent\begin{minipage}{\linewidth}\exercice{8}{Approfondissement}{Famille de tangentes}
{\small\color{Muted}\textbf{Énoncé.} Pour $f(x)=x^3-3x$, chercher les points où la tangente est parallèle à $y=9x+1$.}
\par\medskip
\textbf{Résolution.}\par
Une parallèle a une pente $9$. On résout $f\prime(x)=3x^2-3=9$, soit $x^2=4$, donc $x=-2$ ou $x=2$. Les points sont $(-2,-2)$ et $(2,2)$. Les tangentes valent respectivement $y=9x+16$ et $y=9x-16$.
\end{minipage}\par
\end{document}
